


\def\stat{lukianov}





\def\tit{ПРИКЛАДНЫЕ АСПЕКТЫ КОГНИТИВНОГО МОДЕЛИРОВАНИЯ 


ПРИ~ПРОЕКТИРОВАНИИ СЛОЖНЫХ ИНФОРМАЦИОННЫХ СИСТЕМ}





\def\titkol{Прикладные аспекты когнитивного моделирования 


при~проектировании сложных ИС} %информационных систем}








\def\aut{Г.\,В.~Лукьянов$^1$, Д.\,А.~Никишин$^2$}





\def\autkol{Г.\,В.~Лукьянов, Д.\,А.~Никишин}





\titel{\tit}{\aut}{\autkol}{\titkol}





\index{Лукьянов Г.\,В.}


\index{Никишин Д.\,А.}


\index{Lukyanov G.\,V.}


\index{Nikishin D.\,A.}





%{\renewcommand{\thefootnote}{\fnsymbol{footnote}}


%\footnotetext[1]


%{Работа выполнена при поддержке РФФИ (проект 14-07-00040).}}





\renewcommand{\thefootnote}{\arabic{footnote}}


\footnotetext[1]{Институт проблем информатики Федерального исследовательского


  центра <<Информатика и~управление>> Российской академии наук, 


\mbox{gena-mslu@mail.ru}}


\footnotetext[2]{Институт проблем информатики Федерального исследовательского


  центра <<Информатика и~управление>> Российской академии наук, 


\mbox{dmnikishin@mail.ru}}





\label{st\stat}





                  \thispagestyle{headings}


 


 % \vspace*{-6pt}








\Abst{Статья посвящена прикладным аспектам информационного 


моделирования при проектировании информационных систем (ИС) 


сложных (комплексных) предметных областей. Приведена логическая 


схема проектирования ИС, рассмотрены основные этапы 


проектирования. Представлен гибридный~---  


кас\-кад\-но-ци\-кли\-че\-ский~--- алгоритм формирования мониторинговой 


ИС, призванный компенсировать недостатки как каскадной, так 


и~итерационной парадигмы проектирования и~обеспечить качество 


и~жизнеспособность ИС в~условиях краткосрочного планирования.}





\KW{когнитивное моделирование информационных систем; методы 


моделирования и~разработки информационных систем; условные 


и~безусловные связи}





\DOI{10.14357\!/\!08696527170110} 








\vskip 10pt plus 9pt minus 6pt








      \thispagestyle{headings}


      


    %  \vspace*{-18pt}





\section{Введение}





    Роль моделирования и~моделей в~современной науке и~технике 


заключается в~возможности исследования важнейших свойств 


и~характеристик моделируемой предметной области с~экономически 


обоснованными затратами материальных, финансовых, людских 


и~временн$\acute{\mbox{ы}}$х ресурсов. Независимо от степени слож\-ности предметной 


области моделирование объектов и~процессов предметной области позволяет 


изучать их в~более благоприятной обстановке, например в~лабораторных 


условиях. При этом становится возможным изучение одного и~того же 


объекта или процесса параллельно несколькими исследователями, а~также 


многократно воспроизводить необходимые условия функционирования 


предметной области~[1].


    


    Особую специфику моделирование приобретает при проектировании 


современных ИС, особенно в~тех случаях, 


когда автоматизации в~результате внед\-ре\-ния ИС подлежит слож\-ная, 


многогранная и~мас\-штаб\-ная деятельность. Как показывает опыт создания 


подобных сис\-тем~[1--3], в~этот процесс в~по\-дав\-ля\-ющем большинстве случаев 


приходится вовлекать разнородные, территориально распределенные 


ведомства и~учреждения, которые работают по различным методикам, 


инструкциям и~регламентам. Такая ситуация возникает, например, в~процессе 


автоматизации в~масштабе отдельной отрасли народного хозяйства, когда 


информационный процесс охватывает сотни разнородных показателей, 


которые формируются на основе десятков часто несовместимых методик. 


В~этом и~подобных случаях к~качеству проектирования предъявляются 


повышенные требования, так как последующая переделка недостаточно 


тщательно проработанного проекта ИС связана не только со значительными 


затратами финансовых, материальных и~людских ресурсов, но и~с~потерей 


времени, что является критичным в~системе государственного и~военного 


управ\-ления.


    


    В связи с~этим обеспечение необходимого качества проектирования 


невозможно без тщательно разработанной модели предметной области, что, 


в~свою очередь, требует наличия соответствующей методологии 


моделирования, что и~обусловливает актуальность данной статьи.


    


    Цель исследования состояла в~изучении прикладных аспектов 


моделирования при проектировании сложных ИС, когда одним из решающих 


условий процесса моделирования становится учет когнитивного фактора.


    


    Для достижения сформулированной цели были поставлены следующие 


основные задачи:


    \begin{enumerate}[1.]


    \item Анализ направлений когнитивных исследований в~мировой 


практике.


    \item Раскрытие общих и~частных характеристик модели предметной 


области и~модели ИС.


    \item Выделение наиболее значимых направлений постановки задач при 


проектировании сложных ИС.


    \item Анализ специфики процесса разработки ИС сложных 


и~комплексных предметных областей, преимуществ и~недостатков 


существующих методов проектирования и~формирование практически 


применимого подхода к~проектированию сложных ИС.


    \end{enumerate}





\section{Современные направления когнитивных исследований 


в~структуре~мировой~науки}





    Когнитивная тематика достаточно давно находится в~центре внимания 


ученых и~исследователей~[4--9]. За истекшие~10--15~лет даже появилась 


новая отрасль знаний~--- <<когнитивная информатика>>, которая наиболее 


полно отвечает научным представлениям о взаимосвязи и~взаимозависимости 


мыслительной деятельности человека и~информационного процесса.


    


    Значительная доля исследований посвящена изучению умственных 


механизмов восприятия и~обработки информации и~принятия на этой основе 


решений~[10]. Опора в~научных изысканиях делается, во-пер\-вых, на 


психологические, биологические и~физиологические аспекты деятельности 


человека, что вполне объяснимо. В~качестве яркого примера таких 


исследований может служить работа японских ученых Ф.~Мицогуши, 


Х.~Нишиама и~Х.~Ивасаки <<Новый подход к~обнаружению 


потери внимания водителем на основе данных о движении его глаз>>~[11].


    


    Во-вторых, используются результаты работ в~области искусственного 


интеллекта и~лингвистики в~части структурного анализа текстовой, аудио- 


и~визуальной информации, а также распознавания образов. Этому 


направлению на конференциях, симпозиумах и~семинарах, как правило, 


посвящается отдельная сессия, секция или круглый стол, так как 


исследования здесь носят сильный междисциплинарный характер и~по 


своему содержанию серьезно отличаются от прочих направлений в~области 


когнитивной информатики. Например, в~период формирования форума на 


I~Международной конференции, посвященной когнитивной информатике 


и~когнитивной обработке данных, была проведена специальная сессия 


<<Искусственный интеллект>>~[12]. В~дальнейшем эта традиция 


с~некоторыми вариациями в~названиях неизменно соблюдалась.


    


    Наконец, в-третьих, исследования ведутся в~области 


(автоматизированной) интеллектуальной обработки больших массивов 


неструктурированных данных. Усилия здесь сосредоточены на разработке 


программных средств, имитирующих умственную деятельность человека 


и~позволяющих выявлять неявные функциональные зависимости 


и~неочевидные тенденции на основе анализа больших информационных 


массивов в~интересах разработки объективных и~долгосрочных прогнозов, 


а~также принятия на этой основе адекватных решений. Например, на 13-й 


Международной конференции по когнитивной информатике и~когнитивной 


обработке данных этому направлению было посвящено несколько докладов.


    


    Вместе с~тем ученые и~исследователи пока обходят стороной 


когнитивные аспекты проектирования в~целом и~проектирования и~создания 


ИС в~част\-ности. Хотя именно 


ИС на современном этапе исторического развития служат основным 


инструментом поддержки принятия решений, и~от эффективности их работы 


в~значительной мере зависит качество и~оперативность при\-ни\-ма\-емых 


решений~\cite{2-luk}. Разработка современных сложных 


и~многофункциональных ИС, например для 


мониторинга и~оценки национальной без\-опас\-ности Российской  


Федерации (НБРФ),~--- это многогранный, многоэтапный и~длительный процесс, 


предусматривающий взаимодействие различных структурных подразделений 


исполнителя работ и,~соответственно, многих аналитиков, разработчиков 


и~администраторов, при этом чаще всего функционально разделенных 


и~пространственно распределенных~\cite{1-luk, 3-luk}. И~когнитивный 


аспект в~этой работе играет не последнюю роль.


    


    Таким образом, исследования роли и~значения когнитивной 


составляющей и~ее учета в~процессе проектирования 


ИС даже в~плане формулирования проб\-ле\-мы и~постановки задачи 


является новым, неизученным и~перспективным во-\linebreak просом.





\section{Этапы моделирования}





    Моделирование при проектировании и~создании ИС имеет некоторые 


особенности по сравнению с~другими сферами~[13]. Так, в~самолетостроении 


модель представляет собой образ еще не существующего материального 


объекта~--- самолета. При проектировании же ИС модель прежде всего 


является образом уже существующей сущности (причем как материальной, 


так и~виртуальной) и~отражает эту сущность в~части автоматизации 


происходящих в~ней процессов за счет внедрения ИС. В~данном случае 


модель предметной области служит основой для разработки модели, 


прототипа, проекта будущей ИС.


    


    Исходя из этой специфики и~основываясь на теоретических изысканиях 


в~области применения моделей и~моделирования, логику проектирования, 


разработки и~создания ИС можно представить в~виде последовательности 


действий, изображенной на рис.~1.


    


\begin{figure}[h] %fig1


\vspace*{1pt}


 \begin{center}


 \mbox{%


 \epsfxsize=128.324mm


 \epsfbox{luk-1.eps}


 }


\end{center}


 \vspace*{-9pt}


\Caption{Логическая схема проектирования ИС}


\end{figure}





    Из представленной на рис.~1 логической схемы проектирования ИС 


следует, что по результатам экспериментальных либо теоретических 


исследований предметной области, отраженных в~ее модели, разрабатывается 


модель собственно ИС. Затем проводится проверка соответствия проекта ИС 


целям и~задачам проектирования, специфике предметной области. На 


основании этой проверки уточняется воображаемая модель, а полученные 


изменения воплощаются в~натурной или абстрактной модели с~по\-сле\-ду\-ющей 


корректировкой проекта ИС. 


    


    Описанная выше логика моделирования диктует целесообразность 


сле\-ду\-ющих этапов моделирования~[14, 15].


    


    \textbf{Содержательная постановка задачи.} На этом этапе четко 


формулируются цели и~ставятся задачи моделирования. При этом 


определяются объекты, которые относятся к~решаемой задаче, а также 


ситуация, которую нужно реализовать в~результате ее решения. Для того 


чтобы задачу можно было описать количественно и~использовать при ее 


решении вычислительную технику, нужно произвести качественный 


и~количественный анализ объектов и~ситуаций, имеющих к~ней отношение. 


При этом осуществляется декомпозиция сложных объектов на составные 


части (элементы), определяются взаимосвязи этих элементов, их свойства 


и~количественные и~качественные значения этих свойств, количественные 


и~логические соотношения между ними, выражаемые в~виде уравнений, 


неравенств. В~результате решения этой задачи системного анализа объект 


оказывается представленным в~виде системы.


    


    \textbf{Математическая постановка задачи.} В~процессе 


математической постановки задачи осуществляется построение 


математической модели объекта и~определение методов (алгоритмов) 


получения решения задачи. На этом этапе может оказаться, что ранее 


проведенный системный анализ привел к~такому набору элементов, свойств 


и~соотношений, для которого нет приемлемого метода решения задачи, 


и~в~результате приходится возвращаться к~этапу системного анализа. Как 


правило, решаемые на практике задачи стандартизованы, системный анализ 


производится в~расчете на известную математическую модель и~алгоритм ее 


решения, проблема состоит лишь в~выборе подходящего метода.


    


    \textbf{Разработка программы решения задачи на ЭВМ.} Для 


сложных объектов, состоящих из большого числа элементов и~обладающих 


большим числом свойств, может потребоваться применение 


специализированных пакетов для информационного моделирования, 


обеспечивающих ведение базы данных, средств работы с~ней, методов 


извлечения необходимых для расчетов данных~[16--18]. Для стандартных 


задач может быть использован подходящий пакет прикладных программ 


и~систем управления базами данных.


    


    На заключительном этапе производится опытная эксплуатация модели 


и~на основе анализа полученных результатов реализуется окончательный 


вариант ИС.





\section{Связи между объектами}





    Если проводится глубокая декомпозиция предметной области, 


в~результате которой образуется множество объектов~\cite{17-luk, 18-luk}, 


то отношения между различными классами и~экземплярами объектов 


отражаются в~информационных моделях как связи. Каждая связь задается 


в~модели определенным именем. В~графической форме связь 


представляется как линия со стрелками между связываемыми объектами 


и~обозначается идентификатором связи (например, ER-диа\-грам\-ма, или 


диаграмма <<сущ\-ность--связь>>).


    


    Существуют три вида связи: связь <<один к~одному>> реализуется в~том 


случае, когда один экземпляр одного объекта связан с~единственным 


экземпляром другого; связь <<один ко многим>> возникает тогда, когда один 


экземпляр первого объекта связан с~одним или более экземплярами второго 


объекта, но каждый экземпляр второго объекта связан только с~одним 


экземпляром первого; связь <<многие ко многим>> существует, когда один 


экземпляр первого объекта связан с~одним или несколькими экземплярами 


второго и~каждый экземпляр второго связан с~одним или многими 


экземплярами первого.


    


    Связи подразделяются на безусловные и~условные. В~безусловной связи 


для участия в~ней требуется каждый экземпляр объекта. В~условной связи 


принимают участие не все экземпляры объекта. Связь может быть условной 


как с~одной, так и~с обеих сторон.


    


    Все связи в~информационной модели требуют описания, которое как 


минимум включает:


    \begin{itemize}


\item идентификатор связи;


\item формулировку сущности связи;


\item вид связи (ее множественность и~условность);


\item способ описания связи с~помощью вспомогательных атрибутов 


объектов.


\end{itemize}





    Наиболее известная структура, объединенная однонаправленной 


связью,~--- это очередь. Возможными обобщениями информационных 


моделей являются циклическая структура, таблица, стек.


    


    Очень важную роль играет древовидная информационная модель, 


являющаяся одной из самых распространенных типов классификационных 


структур. Эта модель строится на основе связи, отражающей отношение 


части к~целому, т.\,е.\ связи типа <<один ко многим>>. Таким образом, типы 


данных в~программировании тесно связаны с~определенными 


информационными моделями данных~\cite{16-luk}.


    


    Еще более общей информационной моделью служат так называемые 


графовые структуры, которые лежат в~основе решения многих задач 


информационного моделирования.





\section{Практическая реализация когнитивной составляющей 


моделирования}





    Как было показано в~ряде предыдущих  


работ~\cite{11-luk, 12-luk, 20-luk, 19-luk}, опыт создания сложных ИС, 


предназначенных для решения междисциплинарных задач и~задач 


управления на стратегическом уровне, выявил ряд проблем, затрудняющих 


эффективное проектирование и~разработку таких ИС. Речь прежде всего идет 


о~высоком уровне фрагментарности и~несогласованности связанной с~этой 


деятельностью информации между экспертами, разработчиками 


и~администраторами. В~научной литературе степень такой согласованности 


принято называть уровнем когнитивной интероперабельности~[21--23].


    


    По определению проблему когнитивной интероперабельности позволяет 


преодолеть тщательно разработанная модель предметной  


области~\cite{19-luk}, которая должна предоставить всем участникам 


проекта, независимо от их ведомственной принадлежности 


и~пространственного расположения, целостное и~однозначное представление 


о~самой предметной области, о целях и~задачах проектирования, а~также 


о~функциях, возложенных на ИС. Однако наличие такой модели является 


необходимым, но не достаточным условием. Даже в~случае идеальной 


модели предметной области не удается избавиться от неоднозначного 


понимания тех или иных вопросов проектирования, которые оказывают 


существенное влияние на ход разработки и~на функционирование собственно 


ИС. Это происходит в~силу особенностей персональных знаний экспертов, 


разработчиков и~администраторов, их профессионального опыта. Эти 


особенности восприятия касаются как модели предметной области, так 


и~проекта ИС, что еще больше усугубляет ситуацию.


    


    В качестве эффективного и~наиболее результативного способа 


преодоления этой проблемы предлагается ввести в~регламенты разработки 


модели предметной области и~проектирования ИС раздел, посвященный 


вопросам согласования отдельных частей модели предметной области 


и~проекта. В~этом разделе должны быть предусмотрены следующие 


важнейшие элементы взаимодействия между всеми заинтересованными 


подразделениями и~участниками проектирования:


    \begin{itemize}


\item вопросы, подлежащие обязательному согласованию;


\item порядок изменения (сокращения, расширения) перечня вопросов, 


подлежащих согласованию;


\item формы документов, передаваемых в~порядке согласования между 


участниками проектирования, предусматривающие возможно полное 


и~однозначное описание сути подлежащих согласованию вопросов;


\item жесткие и~конкретные сроки согласования, включая и~возможность 


согласования <<по умолчанию>> в~том случае, если ответ на запрос не 


поступил в~установленный срок;


\item порядок разрешения спорных вопросов в~случае невозможности 


достичь согласованного решения между участниками проектирования.


\end{itemize}





\section{Особенности разработки информационных систем в~целях 


мониторинга и~управления в~сложных предметных областях}


     


     Говоря о создании ИС, важнейшим 


компонентом которой с~точки зрения функциональности является ее 


программное обеспечение (ПО), необходимо отметить ряд широко известных 


методологий (методов и~моделей) процесса разработки 


ПО (\textit{англ.}\ software development process, software process), которые 


в~общем виде можно распространить и~на разработку ИС в~целом, поэтому 


в~дальнейшем речь будет идти о разработке ИС, подразумевая, в~частности, 


разработку соответствующего ПО. 


     


     Индустрия разработки автоматизированных ИС (АИС)


управ\-ле\-ния зародилась в~середине~XX~в., прошла в~своем развитии ряд 


этапов и~приобрела достаточно развитую теоретическую базу для 


организации процессов разработки ИС~\cite{24-luk}. К~началу XXI~в.\ стало 


ясно, что достаточно до\-ро\-го\-сто\-ящие сис\-те\-мы, такие как CAD (Computer 


Aided Design), CAM (Computer Aided Manufacturing), PDM (Product Data 


Management), FRP (Finance Requirements Planning), MRP (Material 


Requirements Planning), MES (Management Execution System), 


ориентированные на локальную автоматизацию и~формирование 


традиционных баз данных, уже не соответствуют требованиям создания 


единого информационного пространства, предназначенного для 


синхронизированного обмена данными всеми участниками 


информационного процесса~\cite{25-luk} и~комплексного использования 


корпоративных данных для целей управления и~планирования 


деятельности~\cite{24-luk}. Таким образом, возникла насущная 


необходимость формирования новой методологии построения~ИС.


     


     В настоящее время ПО АИС предприятия 


разрабатывается на основе концепции PLM (Product Lifecycle Management). 


Анализ возможностей российских и~зарубежных PLM-сис\-тем~\cite{25-luk} 


показывает, что в~них в~той или иной степени реализована интеграция CAD, 


CAM и~тому подобных сис\-тем. В~то же время для осуществления взаимодействия 


в~настоящее время используются методы интеграции на основе бумажной 


технической документации, программного обмена через структурированные 


файлы данных или API (Application Programming Interface). Применение 


таких методов интеграции приводит к~многочисленным ошибкам и~потере 


актуальности данных, существенному затруднению рабочих процессов, 


увеличению стоимости внедрения и~сопровождения. Это обусловливает 


научную проблему, имеющую важное значение для российской 


промышленности в~условиях постоянного роста конкуренции на мировом 


рынке наукоемкой продукции,~--- разработку и~практическую апробацию 


моделей и~методов создания интегрированных АИС, обеспечивающих 


комплексную автоматизацию предприятия~\cite{25-luk}.


     


     Вместе с~тем сложность и~подводные камни процесса проектирования 


объясняют тот факт, что в~настоящее время в~общем потоке проектов 


достаточно мала доля проектов, выполненных в~заявленные сроки 


и~в~рамках заданного объема финансирования. Так, согласно 


статистическим данным, собранным аналитической компанией Standish 


Group (США), из обследованных в~США в~1994~г.\ 8380~проектов, общая 


стоимость которых превышала~80~млрд долларов, неудачными 


оказались более~30\%~проектов~\cite{24-luk}. При этом оказались 


выполненными в~срок лишь~16\% от общего числа проектов, а перерасход 


средств составил~189\% от запланированного бюджета. Из~30\,000~проектов 


АИС, обследованных в~США в~период~1994--2011~гг., успешными 


(выполненными в~срок и~в~рамках заданного объема финансирования) 


оказались не более~24\%. Анализ причин этого показывает~\cite{25-luk}, что 


большинство неудач связано с~отсутствием или неправильным применением 


методологии создания АИС, отвечающей современным требованиям 


предприятий.


     


     Основными проблемами инженерии ПО


становятся ускорение изменений и,~соответственно, необходимость быстрой 


адаптации к~новым условиям, обеспечение, с~одной стороны, высокой 


степени гибкости и~возможности настройки создаваемых систем на быстро 


меняющиеся потребности различных категорий пользователей и~условия 


эксплуатации, а~с~другой~--- функциональной надежности 


разрабатываемых~ИС.


     


     Кроме того, поскольку крупные ИС охватывают все сферы 


деятельности предприятий, необходимым условием создания эффективных 


систем управления, повышения оперативности их адаптации, снижения 


трудоемкости сопровождения является необходимость привлечения 


к~разработке и~модификации моделей специалистов в~различных 


предметных областях.


     


     Также при создании сложных ИС зачастую возникает необходимость 


разработки семейств моделей, описывающих различные стороны 


функционирования сис\-те\-мы и\!/\!или реализации многоуровневых моделей, 


описывающих систему с~разной степенью абстракции~\cite{17-luk}.


     


     Среди множества существующих методологий проектирования ИС 


можно выделить три основных подхода~\cite{24-luk}.


     


     \textbf{Каскадная разработка}, или так называемая <<модель 


водопада>> (\textit{англ.}\ waterfall model), которая выглядит как некий поток работ, 


последовательно проходящий следующие этапы:


     \begin{itemize}


     \item определение требований к~создаваемой ИС;


\item  проектирование будущей ИС, в~ходе которого создается проектная 


документация, описывающая способы и~план реализации определенных 


выше требований;


     \item реализация, т.\,е.\ собственно создание ИС как  


ап\-па\-рат\-но-про\-грам\-мно\-го комплекса;


     \item  тестирование и~отладка ИС~--- на этой стадии в~некоторой 


степени могут быть устранены недочеты, проявившиеся на предыдущих 


стадиях разработки, однако устранение ошибок в~требованиях и~проекте ИС 


представляется весьма проблематичным;


     \item  внедрение ИС и~ее последующая поддержка, осуществляемая 


в~процессе эксплуатации, в~основном направленная на устранение ошибок и~


добавление новой функциональности в~ИС.


     \end{itemize}


    


     Каскадный подход в~чистом виде предполагает, что переход 


к~следующему этапу происходит только после полного и~успешного 


завершения предыдущего, т.\,е.\ возврат к~предыдущим этапам или их 


параллельное выполнение не предусматривается. При управлении большими 


проектами выполнение этапов проекта в~строгой последовательности 


позволяет кардинально снизить многие риски проекта и~сделать более 


прозрачной его реализацию, а~также дает возможность оценивать качество 


продукта на каждом этапе. Тем не менее это обусловливает недостаточную 


гибкость такого подхода и~объявление самоцелью формальное управление 


проектом, зачастую в~ущерб срокам, стоимости или качеству. Для устранения 


этих недостатков был разработан ряд различных модификаций каскадного 


подхода (например, модель водоворота), в~основном заключающихся 


в~обеспечении возможности возврата к~предыдущим этапам, но все они, тем 


не менее, сопряжены со значительным увеличением материальных 


и~временн$\acute{\mbox{ы}}$х затрат на разработку.


     


     \textbf{Спиральная модель} жизненного цикла была предложена для 


преодоления этих проблем. На этапах анализа и~проектирования 


реализуемость технических решений и~степень удовлетворения потребностей 


заказчика проверяется путем создания прототипов. Каждый виток спирали 


соответствует созданию работоспособного фрагмента или версии системы. 


Это позволяет уточнить требования, цели и~характеристики проекта, 


определить качество разработки, спланировать работы следующего витка 


спирали. Таким образом углубляются и~последовательно конкретизируются 


детали проекта и~в~результате выбирается обоснованный вариант, который 


удовлетворяет действительным требованиям заказчика и~доводится до 


реализации. Основная проблема спирального цикла~--- определение момента 


перехода на следующий этап. Для ее решения вводятся временн$\acute{\mbox{ы}}$е 


ограничения на каждый из этапов жизненного цикла, и~переход 


осуществляется в~соответствии с~планом, даже если не вся запланированная 


работа закончена. Планирование производится на основе статистических 


данных, полученных в~предыдущих проектах, и~личного опыта 


разработчиков.


     


     \textbf{Итеративная разработка} (\textit{англ.}\ iteration~--- повторение) 


отражает объективно существующий спиральный цикл создания сложных 


систем и~предполагает выполнение каждого из перечисленных выше этапов 


в~режиме постоянного анализа полученных результатов и~корректировки 


предшествующих этапов работы. Создаваемая таким образом ИС на каждом 


этапе проходит повторяющийся цикл PDCA  


(\textbf{p}lan--\textbf{d}o--\textbf{c}heck--\textbf{a}ct cycle:  


пла\-ни\-ро\-ва\-ние--ре\-а\-ли\-за\-ция--про\-вер\-ка--оцен\-ка). В~ходе 


разработки могут быть изменены существующие или выдвинуты 


дополнительные требования и~ограничения, связанные с~принятыми 


техническими решениями. Именно при итеративной разработке их удается 


учесть в~наиболее полной мере, поскольку именно при таком подходе все 


участники проекта в~полной мере готовы к~изменениям. Разновидностями 


каскадного подхода являются V-мо\-дель, спиральная модель, 


прототипирование.


     


     К преимуществам итеративного подхода также относятся:


     \begin{itemize}


     \item непрерывное итеративное тестирование, раннее обнаружение 


конфликтов между требованиями, моделями и~реализацией проекта, 


снижение воздействия ошибок проектирования на ранних стадиях проекта, 


что ведет к~минимизации затрат на их устранение;


     \item эффективная обратная связь между участниками проекта, 


в~частности проектной команды с~заказчиком ИС, обеспечивающая 


максимально полное удовлетворение его потребностей, эффективное 


использование накопленного опыта;


     \item возможность оценить степень реализации проекта и,~как 


следствие, большая уверенность заказчиков и~исполнителей в~его успешном 


завершении, возможность акцентировать усилия на наиболее важных 


и~критичных направлениях проекта;


     \item более равномерная загрузка участников проекта, более 


равномерное и~предсказуемое распределение затрат по всему периоду 


создания ИС.


     \end{itemize}


     


     Основным недостатком такого подхода является затрудненность 


формального и~поступательного управления проектом, отсутствие в~полной 


мере завершенных стадий практически до самого конца периода разработки, 


что затрудняет детальное распределение бюджета проекта между 


различными этапами и~не позволяет отчитаться о завершении каждого из них 


до окончания всего процесса работы. Это требование не столь существенно 


в~биз\-нес-струк\-ту\-рах, поскольку все работы зачастую могут выполняться 


одним самостоятельным коллективом исполнителей, а заказчик может 


позволить себе не требовать формирования промежуточной отчетности, но 


в~сфере госзаказа, в~условиях жесткого контроля за финансами, это 


требование может стать существенным, а кроме того, в~таких проектах 


зачастую задействовано множество независимых  


кол\-лек\-ти\-вов-ис\-пол\-ни\-те\-лей, что требует четкого распределения 


между ними задач и~финансовых средств.


     


     Все же из-за своих преимуществ итеративный подход сейчас является 


наиболее распространенным. Так, в~стандарте PMBoK (Project Management 


Body of Knowledge), разрабатываемом всемирной некоммерческой 


организацией <<Институт управления проектами>> (Project Management 


Institute, PMI), в~3-й версии формально была закреплена только методика 


<<каскадной модели>>~\cite{26-luk}. Начиная с~PMBoK  


4-й~версии~\cite{27-luk} (вышедшей в~2013~г.)\ предлагаются как 


каскадный, так и~итеративный подходы, а также сочетающая их гибридная 


методология управления проектами. К~тому же, как будет показано ниже, 


разработка ИС для сложных предметных областей практически невозможна 


без применения итеративного подхода.





\section{Гибридный (каскадно-циклический) алгоритм 


формирования информационной системы}





     Обращаясь к~проблематике оценивания и~мониторинга состояния дел 


в~сложных предметных областях (мультидисциплинарных и/или 


междисциплинарных, а также носящих комплексный характер), следует 


отметить, что эти предметные области зачастую далеко не в~полной мере 


изучены или же изучены фрагментарно. Следствием этого является 


отсутствие достаточно полных, но в~то же время общепринятых 


и~практически полезных онтологий (концептуальных схем) 


и~информационных моделей объектов и~процессов, обеспечивающих 


информационное описание этих предметных областей в~создаваемой ИС. 


Более того, возможно также практическое отсутствие очевидных или 


в~достаточной мере проработанных методов и~моделей самого процесса 


мониторинга сложной про\-граммно-це\-ле\-вой деятельности (а~также 


процесса управления ею), учитывающих при\-чин\-но-след\-ст\-вен\-ные 


зависимости и~обеспечивающих правильную постановку целей и~задач 


мониторинга и~управления. Отсутствие таких онтологий, моделей и~методов 


вызывает необходимость применения методов когнитивного моделирования, 


которые, по сути, заложены в~описываемом ниже алгоритме формирования 


проекта~ИС.


     


     В связи со сказанным выше разработка качественной ИС сложной 


предметной области с~ходу, методом каскадного проектирования, является 


делом затруднительным, если не сказать невозможным, и~требует 


применения итерационных методов. С~другой стороны, итерационный 


подход в~чистом виде (описанном выше) носит достаточно общий характер, 


плохо поддается формализации в~плане управления и~бюджетирования. 


Вследствие этого необходимо его конкретизировать применительно как 


к~основной задаче мониторинга сложной предметной области, так 


и~к~возможной задаче управления ею (или ее отдельными субъектами, 


имеющими соответствующие рычаги управления). В~связи с~этим 


предполагается применение гиб\-рид\-но\-го~---  


кас\-кад\-но-цик\-ли\-че\-ско\-го~--- алгоритма формирования мониторинговой 


(или мо\-ни\-то\-рин\-го\-во-управ\-ля\-ющей) ИС для сложных предметных 


областей. Примером такого применения может служить организация 


мониторинга обеспечения НБРФ, в~предметную область которой входят практически все 


без исключения сферы общественной жизни, причем существующие 


в~условиях быстро меняющихся внутренних и~внешних условий, 


приоритетов и~ряда других факторов~\cite{11-luk, 20-luk, 19-luk, 28-luk}.


     


     На рис.~2 представлена блок-схема этого алгоритма, предназначенного 


для мониторинга и~оценивания сложных и~комплексных 


(мультидисциплинарных) предметных областей. 


     


\begin{figure}[b] %fig2


\vspace*{1pt}


 \begin{center}


 \mbox{%


 \epsfxsize=127.496mm


 \epsfbox{luk-2.eps}


 }


\end{center}


 \vspace*{-9pt}


     \Caption{Алгоритм каскадно-циклического процесса формирования ИС 


мониторинга и~управления в~сложной предметной области}


     \end{figure}


     


     Представленные на рис.~2 блоки имеют различную природу:


     \begin{itemize}


     \item[А.] Блок <<Реальная предметная область>> представляет собой 


объекты и~процессы (физические, социальные и~др.), являющиеся предметом 


мониторинга и~управления.\\[-15pt]


     \item[Б.] Блоки <<Цели и~задачи ИС>>, <<Онтология 


и~информационная модель предметной области>> и~<<Алгоритмы 


мониторинга и~управления>>, по сути, представляют собой комплект 


проектной документации, в~котором должно быть формально отражено:\\[-15pt]


     \begin{itemize}


     \item в~первом~--- цели и~задачи процесса мониторинга (и~управления), 


которые одновременно являются целями и~задачами соответствующей ИС;\\[-15pt]


     \item во втором~--- онтология (или концептуальная схема) и~конкретная 


информационная модель объектов и~процессов предметной области, 


реализованная по крайней мере в~той части, которая необходима 


и~достаточна для стоящих целей и~задач;\\[-15pt]


     \item в~третьем~--- соответствующие целям и~задачам алгоритмы 


процессов оценивания и~мониторинга, а~также алгоритмы управляющих 


воздействий, если для ИС предусматривается такая функциональность в~виде 


подсистемы управления.\\[-15pt]


     \end{itemize}


\item[В.] Блок <<Информационная система>> представляет собой  


ап\-па\-рат\-но-про\-граммное воплощение перечисленных выше документов, 


включающее в~себя\footnote{Естественно, следует помнить об организационных 


и~технических (надежность, безопасность и~т.\,п.)\ вопросах создания 


и~функционирования ИС, но в~данном контексте они оставлены за рамками 


рассмотрения.} базу данных, построенную на основе информационной модели, 


и~про\-грам\-мные модули, реализующие алгоритмы мониторинга (и,~при необходимости, 


управления).\\[-15pt]


     \item[Г.] Блок <<Результаты работы>> представляет собой 


генерируемые в~результате практической работы ИС формальные 


документы\footnote{Поскольку речь идет об автоматизированном процессе 


проектирования, то под термином <<документ>> в~первую очередь следует 


понимать документы, существующие и~используемые в~электронном виде 


и~которые лишь при необходимости могут представляться в~бумажном 


виде.}~--- статистические и~аналитические отчеты о~состоянии и~динамике 


предмета мониторинга и~управ\-ления.\\[-15pt]


     \item[Д.] Блок <<Оценка рассогласования>> также представляет собой 


документ (отчет), который отражает степень достижения целей и~задач 


мониторинга и~управления, формально~--- расхождения между требуемыми 


или запланированными значениями и~реально достигнутыми оценками 


показателей~--- результатами мониторинга про\-граммно-це\-ле\-вой 


деятельности. Структура этих расхождений более подробно будет 


рассмотрена ниже\footnote{Вообще говоря, полученные расхождения 


характеризуют не только процесс мониторинга, но и~качество процесса 


управления (если такой функционал предусмотрен в~ИС). Однако если 


механизм управления не предусматривается, то в~сухом остатке остается 


только оценка качества мониторинга.}.


     \end{itemize}


     


     С этими блоками связаны следующие процессы.


     \begin{enumerate}[1.]


     \item Итерационный процесс формирования\footnote[1]{Здесь и~далее под 


термином <<формирование>> понимается, что речь идет не только 


о~создании (проекта, модели, ИС), но и~о~его непрерывной последующей 


модификации.} проекта ИС (тонкие сплошные стрел\-ки). Поскольку речь 


идет о сложных предметных областях, то создание качественной проектной 


документации с~ходу, методом каскадного проектирования, по изложенным 


выше причинам является делом затруднительным, если не сказать 


невозможным. В~связи с~этим предполагается применение итерационного 


циклического подхода уже на данном этапе проектирования в~следующей 


общей последовательности:


     \begin{itemize}


     \item[(а)] предварительно определяются цели и~задачи 


мониторинга\footnote[2]{Опирающиеся на достаточно очевидные 


представления о предмете мониторинга и~управления.};


     \item[(б)] в~соответствии с~целями и~задачами производится изучение 


предметной области на предмет выделения в~ней объектов\!/\!процессов, 


установ\-ле\-ния при\-чин\-но-след\-ст\-вен\-ных связей между объектами 


и~процессами, в~результате чего стро\-ит\-ся/уточ\-ня\-ет\-ся 


информационная модель данной предметной области (определяется 


терминология и~формируется иерархическая структура объектов 


и~процессов);


     \item[(в)] производится анализ полученной информационной модели, на 


основе которого уточняются цели и~задачи ИС (т.\,е.\ осуществляется 


переход вновь к~пункту~(а) и,~при необходимости дальнейшего уточнения, 


к~следующей итерации~--- пунктам~(б) и~(в)).


     \end{itemize}


     Аналогичный процесс выполняется и~для подбора алгоритмов работы 


ИС, однако алгоритмы обычно носят более абстрактный характер, их 


методология достаточно хорошо проработана, поэтому фаза обращения к~


соответствующей предметной области в~общем случае не нужна 


(поэтому~этот цикл показан на рис.~2 одной двусторонней стрелкой).


     \item После достаточно тщательной проработки проекта запускается 


процесс создания ИС (или модификации уже существующей). Этот процесс, 


показанный на рис.~2 толстыми сплошными стрелками, включает разработку 


технической документации, собственно программирование и~развертывание 


ПО на аппаратных платформах, информационное 


наполнение ИС и~сдачу ее в~эксплуатацию.


     \item Основным результатом функционирования мониторинговой ИС 


является периодическое создание информационных отчетов о состоянии 


предмета мониторинга (возможный управляющий функционал ИС в~данной 


работе оставлен за рамками рассмотрения). Процесс создания отчетов 


обозначен на рис.~2 толстой штриховой стрелкой. Основой таких отчетов 


являются таблицы, которые могут дополняться графиками или текстовыми 


разъяснениями. Примером формы такого отчета могут служить 


статистические таблицы в~сборниках Роскомстата.


     \item На основе сформированных мониторинговых отчетов 


формируются вторичные отчеты~--- таблицы рассогласования, отражающие 


различия между запланированными и~реальными значениями показателей 


предмета мониторинга (этот процесс показан на рис.~2 пунктирными 


стрелками). В~общем случае причинами того, что некоторые из результатов 


(значений показателей) оказываются не достигнутыми, могут стать:


     \begin{itemize}


\item неправильная постановка целей и~задач управления;


\item несовершенство системы управления;


\item некорректный подход к~оцениванию или мониторингу результатов.


\end{itemize}


     


     Таким образом, на основе величины рассогласования может быть 


сделан вывод о качестве системы мониторинга, но при этом обязательно 


должно быть исключено влияние несовершенства постановки целей и~задач, 


а также системы управления предметом мониторинга\footnote{Образно 


говоря, можно представить следующую ситуацию: если некий водитель 


имеет целью ускорить движение своего автомобиля, но при этом давит на 


педаль тормоза, а~мониторинг этого процесса осуществляет только по 


показаниям часов, то он может прийти к~иллюзии, что процесс ускорения 


в~принципе идет, хотя и~не с~ожидаемой динамикой (и~усиливает давление 


на тормоз). Этот гротеск призван проиллюстрировать опасность 


формального отношения к~процессу оценивания и~исключительно 


неоднозначную взаимосвязь между перечисленными факторами  


про\-граммно-це\-ле\-вой деятельности.}. Процесс разделения величин 


рассогласования по представленным выше факторам и~оценивание степени 


влияния каждого из них в~настоящее время представляется трудно 


формализуемым, в~связи с~чем его основным инструментом представляется 


метод экспертного оценивания.


     \item На основе анализа рассогласований вырабатываются 


корректирующие воздействия (показаны на рис.~2 штрих\-пунк\-тир\-ны\-ми 


стрелками), оказываемые, во-первых, на формулировку целей и~задач ИС, 


а~во-вто\-рых, на дальнейшее уточнение информационной модели 


и~алгоритмов мониторинга и~оценивания (а~также управления), с~


последующей реализацией всех этих изменений на ап\-па\-рат\-но-про\-граммном 


уровне ИС.


     \end{enumerate}


     


     Представленные блоки и~процессы образуют <<большой>> цикл 


каскадно-цик\-ли\-че\-ско\-го (итерационного) алгоритма формирования ИС, 


предназначенного для мониторинга и~оценивания сложных предметных 


областей. Этот алгоритм был сформирован и~в настоящее время практически 


используется в~процессе разработки макетного образца системы мониторинга и~оценивания в~сфере НБРФ.





\section{Заключение }





    Процесс информационного моделирования должен обеспечивать 


адекватность создаваемой модели целям и~задачам ИС, экономическую 


целесообразность ее реализации, гибкость и~масштабируемость модели 


и~ИС в~целом в~условиях меняющихся целей и~задач 


функционирования. Для этого были рассмотрены традиционный цикл 


информационного моделирования и~специфика моделей, на примере которых 


были проведены параллели между классами и~особенностями моделей 


и~решаемыми на их основе реальными практическими задачами. На этой 


основе намечены подходы к~пониманию сущности и~особенностей 


моделирования при проектировании ИС. Дальнейшее совершенствование 


информационного моделирования связано с~развитием представлений 


о~связях, структурах и~задачах, которые могут быть решены на базе этих 


связей и~структур.


    


    В качестве эффективного и~наиболее результативного способа 


преодоления проблемы неоднозначного понимания различных аспектов 


проектирования, которые оказывают существенное влияние на ход 


разработки и~на функционирование собственно ИС, предлагается ввести 


в~регламенты разработки модели предметной области и~проектирования ИС 


раздел, посвященный вопросам согласования отдельных частей модели 


предметной области и~проекта. 


    


    В работе также было показано, что при создании сложных 


ИС высокого уровня применение в~чистом виде 


традиционных каскадных или циклических методов разработки 


затруднительно. Одной из причин этого является отсутствие или 


недостаточная разработанность соответствующих концептуальных схем 


и~информационных моделей. С~другой стороны, в~условиях  


про\-граммно-це\-ле\-во\-го подхода к~разработке и~созданию подобных ИС, 


особенно при краткосрочном планировании, наиболее подходящим является 


каскадный подход, что не только снижает качество создаваемой ИС, но 


и~прерывает жизненный цикл такой ИС. В~связи с~этим в~статье рассмотрен 


алгоритм кас\-кад\-но-цик\-ли\-че\-ско\-го проектирования, призванный 


устранить перечисленные проблемы и~обеспечить качество 


и~жизнеспособность ИС в~условиях краткосрочного планирования. 


Пред\-став\-лен\-ная схема позволяет понять роль и~место каждого метода 


проектирования, а также целесообразность их применения и~сочетания 


в~различных условиях и~обстоятельствах. Этот алгоритм нашел свое 


практическое применение в~процессе разработки макетного образца системы 


мониторинга и~оценивания в~сфере НБРФ.


    


    Таким образом, сформулированная цель исследования была достигнута, 


а~поставленные задачи решены. Кроме того, в~работе в~схематичном виде 


представлен алгоритм кас\-кад\-но-цик\-ли\-че\-ско\-го проектирования 


сложных ИС, учитывающий когнитивные особенности этого вида 


деятельности. 





{\small\frenchspacing


{%\baselineskip=10.8pt


\addcontentsline{toc}{section}{Литература}


\begin{thebibliography}{99}


\bibitem{1-luk}


\Au{Лукьянов Г.\,В., Никишин Д.\,А.} Когнитивные аспекты информационного 


моделирования при проектировании сложных информационных систем~/\!\!/ 


Системы и~средства информатики, 2016. Т.~26. №\,2. С.~158--170.


\bibitem{2-luk}


\Au{Авдеева З.\,К., Коврига С.\,В., Макаренко~Д.\,И., Максимов~В.\,И.} Когнитивный 


подход в~управлении~/\!\!/ Control Sciences, 2007. №\,3. С.~2--8. 


\bibitem{3-luk}


\Au{Лукьянов Г.\,В., Никишин Д.\,А., Веревкин~Г.\,Ф., Косарик~В.\,В.} Нормативные 


и~методологические аспекты организации информационного мониторинга 


национальной безопасности~/\!\!/ Системы и~средства информатики, 2015. Т.~25. 


№\,3. С.~225--241.


\bibitem{4-luk}


\Au{Axelrod R.} Schema theory: An information processing model of perception and 


cognition~/\!\!/ Am. Political Sci. Rev., 1973. Vol.~67. No.\,4. P.~1248--1266.





\bibitem{9-luk} %5


Structure of decision: The cognitive maps of political elites~/


Ed.\ R.~Axelrod.~--- Princeton, NJ, USA: Princeton University Press, 1976. 404~p.


\bibitem{5-luk} %6


\Au{Axelrod R., Hamilton W.\,D.} The evolution of cooperation~/\!\!/ Science, 1981. 


Vol.~211. No.\,4489. P.~1390--1396.





\bibitem{7-luk}


\Au{Anderson J.\,R., Boyle C.\,F., Corbett~A.\,T., Lewis~M.\,W.} 


Cognitive modeling and intelligent tutoring~/\!\!/ 


Artificial Intelligence, 1990. Vol.~42. No.\,1. P.~7--49.





\bibitem{6-luk} %8


\Au{Funge J., Tu X., Terzopoulos~D.} Cognitive modeling: Knowledge, reasoning and 


planning for intelligent characters~/\!\!/ 26th Annual Conference on Computer Graphics 


and Interactive Techniques Proceedings.~--- ACM Press/Addison-Wesley, 1999.  


P.~29--38.





\bibitem{8-luk} %9


\Au{Jakobson G., Buford J., Lewis~L.} A~framework of cognitive situation modeling and 


recognition~/\!\!/ IEEE Military Communications Conference (MILCOM-2006).~--- IEEE, 


2006. P.~1--7.





\bibitem{10-luk}


\Au{Tsvetkov V.\,Ya.} The cognitive modeling with the use of spatial information~/\!\!/ 


Eur. J.~Technology Design, 2015. Vol.~10. Iss.~4. P.~149--158.


\bibitem{11-luk}


\Au{Mizoguchi F., Nishiyama H., Iwasaki~H.} A~new approach to detecting distracted car 


drivers using eye-movement data~/\!\!/ 13th IEEE  Conference (International) on 


Cognitive Informatics and Cognitive Computing Proceedings.~--- London, U.K.:


IEEE, 2014. P.~266--272.


\bibitem{12-luk}


Session 3(B): Artificial Intelligence (AI). 1st IEEE Conference (International) on 


Cognitive Informatics Proceedings.~--- Calgary, Alberth, Canada:


IEEE, 2002. P.~295--328.


\bibitem{13-luk}


\Au{Городецкий А.\,Я.} Информационные системы. Вероятностные модели 


и~статистические решения.~--- СПб.: СПбГПУ, 2003. 326~c.


\bibitem{15-luk} %14


\Au{Ипатова~Э.\,Р.} Методологии и~технологии системного проектирования 


информационных сис\-тем.~--- М.: ФЛИНТА, 2008. 256~с.


\bibitem{14-luk} %15


\Au{Аверченков~В.\,И.} Основы математического моделирования технических  


сис\-тем.~--- М.: ФЛИНТА, 2011. 271~с.








\bibitem{17-luk} %16


\Au{Лядова~Л.\,Н., Сухов~ А\, О.} Визуальные языки и~языковые инструментарии: 


методы и~средства реализации~/\!\!/ Тр. Конгресса по интеллектуальным системам 


и~информационным технологиям AIS-IT, 2010. Т.~10. С.~374--382.





\bibitem{16-luk} %17


\Au{Петеляк~В.\,Е., Новикова~Т.\,Б., Масленникова~О.\,Е., Махмутова~М.\,В., 


Аг\-дав\-ле\-то\-ва~А.\,М.} Data flow diagramming: Особенности построения моделей 


описания управления потоками данных в~организационных сис\-те\-мах~/\!\!/ 


Фундаментальные исследования, 2015. №\,8. С.~323--327.





\bibitem{18-luk}


\Au{Замятина~Е.\,Б., Лядова~Л.\,Н., Сухов~А.\,О.} О~подходе к~интеграции систем 


моделирования и~информационных систем на основе DSM-плат\-фор\-мы 


MetaLanguage~/\!\!/ New Voices Translation Studies, 2016. №\,14. С.~46--73.


\bibitem{20-luk} %19


\Au{Лукьянов~Г.\,В., Никишин~Д.\,А., Веревкин~Г.\,Ф., Косарик~В.\,В.} Специфика 


показателей национальной безопасности в~контексте ее информационного 


мониторинга~/\!\!/ Системы и~средства информатики, 2014. Т.~24. №\,4. С.~186--205.





\bibitem{19-luk} %20


\Au{Лукьянов~Г.\,В., Никишин~Д.\,А., Веревкин~Г.\,Ф., Косарик~В.\,В.} Проб\-ле\-ма 


нормализации показателей информационного мониторинга национальной  


без\-опас\-ности~/\!\!/ Системы и~средства информатики, 2015. Т.~25. №\,4.  


С.~201--213.


\bibitem{21-luk}


\Au{Buddenberg~R.} Toward an interoperability reference model. 2006. {\sf 


http://www.\linebreak dtic.mil/cgi-bin/GetTRDoc?AD=ADA502999}.


\bibitem{22-luk}


\Au{Goldkuhl~G.} The challenges of interoperability in E-government: Towards 


a~conceptual refinement~/\!\!/ Pre-ICIS SIG eGovernment Workshop Proceedings,  2008. 


{\sf 


https:// www.researchgate.net/profile/Goeran\_Goldkuhl/publication/228903679\_The\_\linebreak


challenges\_of\_Interoperability\_in\_E-government\_Towards\_a\_conceptual\_refinement/\linebreak 


links/5592907908ae1e9cb4295340.pdf}.


\bibitem{23-luk}


\Au{Дулин~С.\,К., Дулина Н.\,Г., Кожунова~О.\,С.} Когнитивная 


интероперабельность экспертной деятельности и~ее приложение 


в~геоинформатике~/\!\!/ Тр. 13-й конф. по искуственному  


интеллекту.~--- Белгород: БГТУ, 2012. Т.~1. С.~351--357.


\bibitem{24-luk}


\Au{Грекул~В.\,И., Денищенко~Г.\,Н., Коровкина~Н.\,Л.} Проектирование 


информационных сис\-тем.~--- М.: БИНОМ, 2008. 304~с.


\bibitem{25-luk}


\Au{Кульга~К.\,С.} Модели и~методы создания интегрированной информационной 


сис\-те\-мы для автоматизации технической подготовки и~управ\-ле\-ния 


авиационным производством~/\!\!/ Известия Самарского научного центра РАН, 2012. 


Т.~14. №\,4(2). С.~437--445.


\bibitem{26-luk}


\Au{Whitty S.\,J., Schulz~M.\,F.} THE\_PM\_BOK\_CODE~/\!\!/ 20th IPMA World 


Congress on Project Management, 2006. Vol.~1. P.~466--472.


\bibitem{27-luk}


A~guide to the project management body of knowledge (PMBOK guide).~--- 5th ed.~--- 


Project Management Institute, 2013. 589~p.


\bibitem{28-luk}


\Au{Замятина~Е.\,Б., Лядова~Л.\,Н., Сухов~А.\,О.} Интеграция систем 


моделирования на основе DSM-плат\-фор\-мы с~использованием онтологий~/\!\!/ 


Открытые семантические технологии проектирования интеллектуальных сис\-тем 


(OSTIS-2014): Мат-лы IV~Междунар. научн.-технич. конф.~--- Минск: БГУИР, 2014. 


С.~375--380.


       \end{thebibliography}


} }








\vspace*{-3pt}





\hfill{\small\textit{Поступила в~редакцию 11.07.16}}














\vspace*{9pt}





\hrule





%\newpage





\vspace*{2pt}





\hrule





%\vspace*{9pt}











\def\tit{APPLIED ASPECTS OF~MODELING OF~INFORMATION SYSTEMS}





\def\titkol{Applied aspects of~modeling of~information systems}








\def\aut{G.\,V.~Lukyanov and~D.\,A.~Nikishin}


\def\autkol{G.\,V.~Lukyanov and~D.\,A.~Nikishin}











\titel{\tit}{\aut}{\autkol}{\titkol}





\vspace*{-9pt}








\noindent


Institute of Informatics Problems, Federal Research Center ``Computer Science


and Control'' of the Russian Academy of 


Sciences, 44-2~Vavilov Str., Moscow 119333, Russian Federation











\def\leftfootline{\small{\textbf{\thepage}


\hfill Sistemy i Sredstva Informatiki~---


Systems and Means of Informatics 2017 vol~27 no\,1}


}%


 \def\rightfootline{\small{Sistemy i Sredstva Informatiki~---


 Systems and Means of Informatics   2017   vol~27   no\,1


\hfill \textbf{\thepage}}}











\Abste{The article is devoted to applied aspects of information modeling in the 


design of information systems of complex subject areas. It provides a logic 


diagram of the information system design and substantial description of the basic 


design stages. The article presents the hybrid cascaded-round-Robin algorithm of 


formation of a monitoring (monitoring, command, and control) information system. 


This algorithm is designed to compensate for the shortcomings of the cascade and 


iterative paradigms of design and to ensure the quality and viability of an 


information system in terms of short-term planning.}





\KWE{modeling of information systems; iterative development; waterfall 


development; conditional and unconditional links}





\DOI{10.14357\!/\!08696527170110} 





%\vspace*{-24pt}





% \Ack


% \noindent





\renewcommand{\bibname}{\protect\rmfamily References}


%\renewcommand{\bibname}{\large\protect\rm References}





{\small\frenchspacing


{%\baselineskip=10.8pt


\addcontentsline{toc}{section}{References}


\begin{thebibliography}{99}


\bibitem{1-luk-1}


\Aue{Luk'yanov,~G.\,V., and D.\,A.~Nikishin}. 2015. Kognitivnye 


aspekty informatsionnogo modelirovaniya pri proektirovanii slozhnykh 


informatsionnykh sistem [Cognitive aspects of information modeling in 


the design of complex information systems]. \textit{Sistemy i~Sredstva 


Informatiki~--- Systems and Means of Informatics} 26(2):158--170.


\bibitem{2-luk-1}


\Aue{Avdeeva, Z.\,K., S.\,V.~Kovriga, D.\,I.~Makarenko, and V.\,I.~Maksimov}. 


2007. Kognitivnyy podkhod v~upravlenii [Cognitive approach in 


management]. \textit{Control Sciences} 3:2--8.


\bibitem{3-luk-1}


\Aue{Luk'yanov, G.\,V., D.\,A.~Nikishin, G.\,F.~Verevkin, and V.\,V.~Kosarik}. 


2015. Normativnye i~metodologicheskie aspekty organizatsii 


informatsionnogo monitoringa na\-tsi\-o\-nal'\-noy bezopasnosti [Standard and 


methodological aspects of the organization of information monitoring of 


national security]. \textit{Sistemy i~Sredstva Informatiki~--- Systems and 


Means of Informatics} 25(3):225--241.


\bibitem{4-luk-1}


\Aue{Axelrod, R.} 1973. Schema theory: An information processing model of 


perception and cognition. \textit{Am. Political Sci. Rev.}  


\bibitem{9-luk-1} %5


Axelrod, R., ed. 1976. \textit{Structure of decision: The cognitive maps of political 


elites}. Princeton, NJ: Princeton University Press. 404~p.


67(4):1248--1266.


\bibitem{5-luk-1} %6


\Aue{Axelrod, R., and W.\,D.~Hamilton}. 1981. The evolution of cooperation. 


\textit{Science} 211(4489):1390--1396.


\bibitem{7-luk-1} %7


\Aue{Anderson, J.\,R., C.\,F.~Boyle, A.\,T.~Corbett, and M.\,W.~Lewis.} 


1990. Cognitive modeling and intelligent 


tutoring. \textit{Artificial Intelligence} 42(1):7--49.


\bibitem{6-luk-1} %8


\Aue{Funge, J., X.~Tu, and D.~Terzopoulos}. 1999. Cognitive modeling: 


Knowledge, reasoning and planning for intelligent characters. \textit{26th 


Annual Conference on Computer Graphics and Interactive Techniques 


Proceedings}. ACM Press/Addison-Wesley Publishing Co. 29--38.





\bibitem{8-luk-1} %9


\Aue{Jakobson, G., J. Buford, and L.~Lewis.} 2006. A~framework of cognitive 


situation modeling and recognition. \textit{IEEE Military Communications 


Conference}.  IEEE. 1--7.





\bibitem{10-luk-1}


\Aue{Tsvetkov, V.\,Ya.} 2015. The cognitive modeling with the use of spatial 


information. \textit{Eur.~J. Technology Design} 10(4):149--158.


\bibitem{11-luk-1}


\Aue{Mizoguchi, F., H.~Nishiyama, and H.~Iwasaki.} 2014. A~new approach 


to detecting distracted car drivers using eye-movement data. \textit{13th 


IEEE  Conference (International) on Cognitive Informatics and Cognitive 


Computing Proceedings}. London, U.K.: IEEE. 266--272.


\bibitem{12-luk-1}


Session 3(B): Artificial Intelligence (AI). 2002.


\textit{First IEEE International 


Conference on Cognitive Informatics Proceedings}.  Calgary, Alberta, Canada: 


IEEE. 295--328.


\bibitem{13-luk-1}


\Aue{Gorodetskiy, A.\,Ya.} 2003. \textit{Informatsionnye sistemy. 


Veroyatnostnye modeli i~sta\-ti\-sti\-che\-skie resheniya} [Information systems. 


Probabilistic models and statistical decisions]. St.\ Petersburg: SPbGPU


Publs. 326~p. 





\bibitem{15-luk-1} %14


\Aue{Ipatova, E.\,R.} 2008. \textit{Metodologii i~tekhnologii sistemnogo 


proektirovaniya in\-for\-ma\-tsi\-on\-nykh sistem} [Methodology and technology of 


system design of information systems]. Moscow: FLINTA. 256~p.


\bibitem{14-luk-1} %15


\Aue{Averchenkov, V.\,I.} 2011. \textit{Osnovy matematicheskogo 


modelirovaniya tekhnicheskikh sistem} [Foundations of mathematical modeling 


of technical systems]. Moscow: FLINTA. 271~p.








\bibitem{17-luk-1} %16


\Aue{Lyadova, L.\,N., and A.\,O.~Sukhov}. 2010. Vizual'nye yazyki 


i~yazykovye instrumentarii: metody i~sredstva realizatsii [Visual languages 


and language tools: Methods and tools for implementation]. \textit{Tr. 


Kongressa po intellektual'nym sistemam i~informatsionnym tekhnologiyam 


AIS-IT} [Congress on Intelligent Systems and Information Technologies AIS-IT 


Proceedings]. 10:374--382.


\bibitem{16-luk-1} %17


\Aue{Petelyak, V.\,E., T.\,B.~Novikova, O.\,E.~Maslennikova, 


M.\,V.~Makhmutova, and A.\,M.~Agdavletova}. 2015. Data flow 


diagramming: Osobennosti postroeniya mo\-de\-ley opisaniya upravleniya 


potokami dannykh v~organizatsionnykh sistemakh [Data flow diagramming: 


Features of building models for the description of data flow management in 


organizational systems]. \textit{Fundametnal'nye issledovaniya} [Fundamental 


Researsh] 8:323--327.





\bibitem{18-luk-1} %18


\Aue{Zamyatina, E.\,B., L.\,N.~Lyadova, and A.\,O.~Sukhov}. 2016. 


O~podkhode k~integratsii sistem modelirovaniya i~informatsionnykh sistem na 


osnove DSM-platformy MetaLanguage [On the approach to the integration of 


systems modeling and information systems on the basis of DSM-platform 


MetaLanguage]. \textit{New Voices Translation Studies} 14:46--73.





\bibitem{20-luk-1} %19


\Aue{Luk'yanov, G.\,V., D.\,A. Nikishin, G.\,F.~Verevkin, and 


V.\,V.~Kosarik}. 2014. Spetsifika pokazateley natsional'noy bezopasnosti 


v~kontekste ee informatsionnogo monitoringa [Specificity of indicators of 


national security in the context of its information monitoring]. \textit{Sistemy 


i~Sredstva Informatiki~--- Systems and Means of Informatics} 24(4):186--205.


\bibitem{19-luk-1} %20


\Aue{Luk'yanov, G.\,V., D.\,A. Nikishin, G.\,F.~Verevkin, and 


V.\,V.~Kosarik}. 2015. Problema normalizatsii pokazateley informatsionnogo 


monitoringa natsional'noy bezopasnosti [Problem of normalization of indicators 


of information monitoring of national security]. \textit{Sistemy i~Sredstva 


Informatiki~--- Systems and Means of Informatics} 25(4):201--213.





\bibitem{21-luk-1}


\Aue{Buddenberg, R.} 2006. Toward an interoperability reference model. 


Available at:  


{\sf http:// www.dtic.mil/cgi-bin/GetTRDoc?AD=ADA502999} (accessed 


September~11, 2016).


\bibitem{22-luk-1}


\Aue{Goldkuhl, G.} 2008. The challenges of interoperability in E-government: 


Towards a~conceptual refinement. \textit{Pre-ICIS SIG eGovernment 


Workshop Proceedings}. Available at: {\sf  


https://www.researchgate.net/profile/Goeran\_Goldkuhl/publication/228903679\_\linebreak 


The\_challenges\_of\_Interoperability\_in\_E-government\_Towards\_a\_conceptual\_\linebreak


refinement/links/5592907908ae1e9cb4295340.pdf} 


(accessed September~11, 2016).


\bibitem{23-luk-1}


\Aue{Dulin, S.K\,., N.\,G. Dulina, and O.\,S.~Kozhunova}. 2012. Kognitivnaya 


in\-ter\-ope\-ra\-bel'\-nost' ekspertnoy deyatel'nosti i~ee prilozhenie v~geoinformatike 


[Cognitive interoperability expert activity and its application in 


geoinformatics]. \textit{Tr. 13-y konf. po iskustvennomu intellektu} [13th 


Conference on Artificially Intelligence Proceedings].


Belgorod: BFTU Publs. 1:351--357.


\bibitem{24-luk-1}


\Aue{Grekul, V.\,I., G.\,N. Denishchenko, and N.\,L.~Korovkina}. 2008. 


\textit{Proektirovanie informatsionnykh sistem} [Design of information 


systems]. Moscow: Binom. 304~p.


\bibitem{25-luk-1}


\Aue{Kul'ga, K.\,S.} 2012. Modeli i~metody sozdaniya integrirovannoy 


informatsionnoy sistemy dlya avtomatizatsii tekhnicheskoy podgotovki 


i~upravleniya aviatsionnym proizvodstvom [Models and methods of creation 


of integrated information system for automation of technical training and 


control of aviation production]. \textit{Izv. Samarskogo nauchnogo tsentra 


RAN} [Izv. Samara Scientific Center of RAS] 14(2):437--445.


\bibitem{26-luk-1}


\Aue{Whitty, S.\,J., and M.\,F.~Schulz}. 2006. THE\_PM\_BOK\_CODE. 


\textit{20th IPMA World Congress on Project Management}. 1:466--472.


\bibitem{27-luk-1}


Project Management Institute. 2013.


 \textit{A~guide to the project management body of knowledge (PMBOK 


guide)}. 5th ed.  589~p.


\bibitem{28-luk-1}


\Aue{Zamyatina, E.\,B., L.\,N.~Lyadova, and A.\,O.~Sukhov}. 2014. 


Integratsiya sistem modelirovaniya na osnove DSM-platformy s ispol'zovaniem 


ontologiy [Integration of modeling systems on the basis of DSM-platform with 


ontologies]. \textit{Conference (International)


``Open Semantic Technologies for Intelligent Systems'' 


Proceedings}. Minsk: BGUIR. 375--380.


 \end{thebibliography}


} }








\vspace*{-3pt}





\hfill{\small\textit{Received July 11, 2016}}   


    


\Contr





\noindent


\textbf{Lukyanov Gennady V.} (b.\ 1952)~--- Candidate of Military Science (PhD), 


associate professor, leading scientist, Institute of Informatics Problems, Federal 


Research Center ``Computer Science and Control'' of the Russian Academy of 


Sciences, 44-2~Vavilov Str., Moscow 119333, Russian Federation; 


\mbox{gena-mslu@mail.ru}





\vspace*{3pt}





\noindent


\textbf{Nikishin Dmitry A.} (b.\ 1976)~--- Candidate of Science (PhD) in technology, 


leading scientist, Institute of Informatics Problems, Federal Research Center ``Computer Science and 


Control'' of the Russian Academy of Sciences, 44-2~Vavilov Str., Moscow 119333, Russian 


Federation; \mbox{dmnikishin@mail.ru}





\label{end\stat}








\renewcommand{\bibname}{\protect\rm Литература}    


